
\documentclass{article}
\usepackage{colortbl}
\usepackage{makecell}
\usepackage{multirow}
\usepackage{supertabular}

\begin{document}

\newcounter{utterance}

\centering \large Interaction Transcript for game `cladder', experiment `full\_v1.5\_default', episode 7275 with Qwen3.6{-}27B.
\vspace{24pt}

{ \footnotesize  \setcounter{utterance}{1}
\setlength{\tabcolsep}{0pt}
\begin{supertabular}{c@{$\;$}|p{.15\linewidth}@{}p{.15\linewidth}p{.15\linewidth}p{.15\linewidth}p{.15\linewidth}p{.15\linewidth}}
    \# & \multicolumn{2}{c}{Player} && \multicolumn{2}{c}{Game Master} \\
    \hline

    \theutterance \stepcounter{utterance}  
    & & & \multicolumn{4}{p{0.6\linewidth}}{
        \cellcolor[rgb]{0.9,0.9,0.9}{
            \makecell[{{p{\linewidth}}}]{
                \texttt{\tiny{[P1$\langle$GM]}}
                \texttt{You are an expert in causal inference. The following question is not a typical commonsense query, but rather a meticulously designed question created by a professor specializing in causal inference, intended to assess the students' mastery of the course content.} \\
\texttt{Imagine a self{-}contained, hypothetical world with only the following conditions, and without any unmentioned factors or causal relationships: Yupt has a direct effect on zupj and muvq. Zupj has a direct effect on muvq. The overall probability of yupt is 1\%. For those who are not yupt, the probability of muvq is 33\%. For those who are yupt, the probability of muvq is 59\%. Is muvq more likely than not muvq overall?} \\
\texttt{Start your answer with 'Yes' or 'No', followed by additional reasoning or evidence to support your explanation.} \\
            }
        }
    }
     \\ \\

    \theutterance \stepcounter{utterance}  
    & \multicolumn{4}{p{0.6\linewidth}}{
        \cellcolor[rgb]{0.561,0.737,0.561}{
            \makecell[{{p{\linewidth}}}]{
                \texttt{\tiny{[P1$\rangle$GM]}}
                \texttt{No} \\
\\ 
\texttt{To determine whether `muvq` is more likely than not `muvq` overall, we need to calculate the total probability of `muvq`, denoted as \$P(\textbackslash{}text\{muvq\})\$.} \\
\\ 
\texttt{We are given the following information:} \\
\texttt{1.  \$P(\textbackslash{}text\{yupt\}) = 0.01\$} \\
\texttt{2.  Therefore, \$P(\textbackslash{}neg \textbackslash{}text\{yupt\}) = 1 {-} 0.01 = 0.99\$} \\
\texttt{3.  \$P(\textbackslash{}text\{muvq\} | \textbackslash{}neg \textbackslash{}text\{yupt\}) = 0.33\$} \\
\texttt{4.  \$P(\textbackslash{}text\{muvq\} | \textbackslash{}text\{yupt\}) = 0.59\$} \\
\\ 
\texttt{Using the Law of Total Probability, we can calculate \$P(\textbackslash{}text\{muvq\})\$ as follows:} \\
\\ 
\texttt{\$\$} \\
\texttt{P(\textbackslash{}text\{muvq\}) = P(\textbackslash{}text\{muvq\} | \textbackslash{}text\{yupt\}) \textbackslash{}cdot P(\textbackslash{}text\{yupt\}) + P(\textbackslash{}text\{muvq\} | \textbackslash{}neg \textbackslash{}text\{yupt\}) \textbackslash{}cdot P(\textbackslash{}neg \textbackslash{}text\{yupt\})} \\
\texttt{\$\$} \\
\\ 
\texttt{Substituting the given values:} \\
\\ 
\texttt{\$\$} \\
\texttt{P(\textbackslash{}text\{muvq\}) = (0.59 \textbackslash{}cdot 0.01) + (0.33 \textbackslash{}cdot 0.99} \\
            }
        }
    }
    & & \\ \\

    \theutterance \stepcounter{utterance}  
    & & & \multicolumn{2}{p{0.3\linewidth}}{
        \cellcolor[rgb]{0.9,0.9,0.9}{
            \makecell[{{p{\linewidth}}}]{
                \texttt{\tiny{[GM$|$GM]}}
                \texttt{No} \\
            }
        }
    }
    & & \\ \\

    \theutterance \stepcounter{utterance}  
    & & & \multicolumn{2}{p{0.3\linewidth}}{
        \cellcolor[rgb]{0.9,0.9,0.9}{
            \makecell[{{p{\linewidth}}}]{
                \texttt{\tiny{[GM$|$GM]}}
                \texttt{no} \\
            }
        }
    }
    & & \\ \\

    \theutterance \stepcounter{utterance}  
    & & & \multicolumn{2}{p{0.3\linewidth}}{
        \cellcolor[rgb]{0.9,0.9,0.9}{
            \makecell[{{p{\linewidth}}}]{
                \texttt{\tiny{[GM$|$GM]}}
                \texttt{game\_result = WIN} \\
            }
        }
    }
    & & \\ \\

\end{supertabular}
}

\end{document}
